%=======================================================================
%  Slides 3-15 - Bài toán Downstream Task
%   Transcribed 1:1 from slide-cs222.pptx.pdf (same order, same wording).
%=======================================================================
\section{Bài toán Downstream Task}

% ------------------------------------------------------------------ 3
\begin{frame}{Nội dung chương 1}
  \begin{columns}[T,onlytextwidth]
    \begin{column}{0.5\textwidth}
      \begin{enumerate}\tightlist
        \item[\textbf{1.1}] Giới thiệu bài toán
        \item[\textbf{1.2}] Phương pháp
        \item[\textbf{1.3}] Hạn chế
      \end{enumerate}
    \end{column}
    \begin{column}{0.44\textwidth}
      \begin{beamercolorbox}[wd=\dimexpr\linewidth-1em\relax,sep=0.6em]{block body}
        \footnotesize\color{Muted}
        \centering 1. Bài toán Downstream Task
      \end{beamercolorbox}
    \end{column}
  \end{columns}
\end{frame}

% ------------------------------------------------------------------ 4
\begin{frame}{1.1\quad Giới thiệu bài toán}
  \begin{columns}[T,onlytextwidth]
    \begin{column}{0.53\textwidth}
      \justifying
      Model là một mô hình ngôn ngữ đã qua ký hiệu là $M$, sở hữu bộ trọng số gốc
      $\theta_0$.

      \vspace{0.35em}
      Đầu ra của model $M$ trên 1 đầu vào $x$ bị ảnh hưởng bởi 2 thành phần:
      \begin{itemize}\tightlist
        \item $\pi$: Prompt (chỉ dẫn hay ngữ cảnh ở không gian đầu vào)
        \item denta\_$\theta$: Phần cập nhật trọng số.
      \end{itemize}
    \end{column}
    \begin{column}{0.43\textwidth}
      \begin{beamercolorbox}[wd=\dimexpr\linewidth-1em\relax,sep=0.7em]{block body}
        \centering\footnotesize
        Đầu ra của model $M$ với input $x$:\\[0.4em]
        {\large $y = M(x;\ \pi,\ \theta_0 + \Delta\theta)$}
      \end{beamercolorbox}
    \end{column}
  \end{columns}
\end{frame}

% ------------------------------------------------------------------ 5
\begin{frame}{1.1\quad Giới thiệu bài toán}
  \justifying
  \textbf{Downsteam Task Adaptation} là quá trình điều chỉnh một Model ($M$) để nó
  giải quyết tốt một bài toán cụ thể ($T$), thay vì phải huấn luyện một mô hình mới
  từ đầu.

  \vspace{0.3em}
  Đặt $D_T$ là tập phân phối dữ liệu $(x, m)$ của task $T$.
  \begin{itemize}\tightlist
    \item Với $x$ là input của task $T$.
    \item Với $m$ là đáp án đúng của task $T$.
  \end{itemize}

  \begin{beamercolorbox}[wd=\dimexpr\linewidth-1em\relax,sep=0.5em]{block body}
    \centering\footnotesize
    Nhiệm vụ:\\[0.3em]
    $(\pi^*, \Delta\theta^*) = \arg\max_{(\pi, \Delta\theta)}\ \mathbb{E}_{((x,\,m)\sim D_T)}\big[\ \mu(M(x;\ \pi,\ \theta_0 + \Delta\theta),\ m)\ \big]$
  \end{beamercolorbox}
  \vspace{0.2em}
  $\Leftrightarrow$ \textbf{Mục tiêu:} tìm bộ trọng số $\theta^*$ và bộ prompt $\pi^*$ tốt nhất.
\end{frame}

% ------------------------------------------------------------------ 6
\begin{frame}{1.1\quad Giới thiệu bài toán}
  Ví dụ một số downstream task:
  \begin{itemize}\tightlist
    \item Suy luận giải Toán
    \item Hỏi đáp suy luận đa bước
    \item Kiểm chứng thông tin
    \item Tuân thủ hướng dẫn phức tạp
    \item Tối ưu hoá mã nguồn
    \item \dots
  \end{itemize}
\end{frame}

% ------------------------------------------------------------------ 7
\begin{frame}{1.1\quad Giới thiệu bài toán: Downstream Task}
\end{frame}

% ------------------------------------------------------------------ 8
\begin{frame}{1.1\quad Giới thiệu bài toán}
  \FigureFull[0.66]{t2l/fig4_qualitative}{Hình 4 --- Text-to-LoRA (ICML 2025)}
\end{frame}

% ------------------------------------------------------------------ 9
\begin{frame}{1.2\quad Phương pháp cho DTA}
  \textbf{Hướng tối ưu prompt ($\pi$)}
  \begin{itemize}\tightlist
    \item In-Context Learning:
      \begin{itemize}\tightlist
        \item Few-shot
        \item Chain of Thought
        \item Self-Consistency
        \item Tree-of-Thoughts
        \item \dots
      \end{itemize}
  \end{itemize}
\end{frame}

% ------------------------------------------------------------------ 10
\begin{frame}{1.2\quad Phương pháp cho DTA}
  \textbf{Hướng tối ưu prompt:}
  \begin{itemize}\tightlist
    \item In-Context Learning:
    \item Automated Optimization:
      \begin{itemize}\tightlist
        \item LLM-as-an-Optimizer: Dùng LLM để tối ưu prompt
        \item EvoPrompt: Dùng thuật toán tiến hoá.
        \item RL-based Prompting: Dùng RL huấn luyện một mô hình nhỏ để sinh ra chuỗi
              token prompt giúp tối đa hóa phần thưởng.
        \item Hệ thống có nhiều LLM (Compound AI System):
          \begin{itemize}\tightlist
            \item Tối ưu đồng bộ cả câu chỉ dẫn lẫn bộ mẫu few-shot cho toàn bộ
                  pipeline gồm nhiều module LLM nối tiếp nhau. (DSPy, MIPROv2).
          \end{itemize}
      \end{itemize}
  \end{itemize}
\end{frame}

% ------------------------------------------------------------------ 11
\begin{frame}{1.2\quad Phương pháp cho DTA}
  \textbf{Hướng tối ưu tham số mô hình ($\theta$)}
  \begin{itemize}\tightlist
    \item Full Fine-Tuning (FFT)
    \item Parameter-Efficient Fine-Tuning (PEFT): Đóng băng gần như toàn bộ mô hình gốc,
          chỉ huấn luyện số lượng tham số $\ll$ tổng số tham số.
      \begin{itemize}\tightlist
        \item Reparameterization: LoRa, QLoRa, \dots
        \item Adapters
        \item Prefix-Tuning / Prompt-Tuning
      \end{itemize}
  \end{itemize}
\end{frame}

% ------------------------------------------------------------------ 12
\begin{frame}{1.3\quad Hạn chế của từng phương pháp}{Hạn chế của hướng Prompt}
  \begin{itemize}\tightlist
    \item In Context Learning thủ công:
      \begin{itemize}\tightlist
        \item Phụ thuộc vào kinh nghiệm thử-sai, cách diễn đạt câu chữ và thứ tự sắp
              xếp ví dụ.
        \item Không biết đưa ví dụ few-shot nào phù hợp.
        \item Hệ thống gồm nhiều LLMs (Compound AI System) thì sẽ rất khó để chỉnh tay
              tất cả.
      \end{itemize}
  \end{itemize}
  \vspace{0.3em}
  $\rightarrow$ Automated Optimization ra đời.
\end{frame}

% ------------------------------------------------------------------ 13
\begin{frame}{1.3\quad Hạn chế của từng phương pháp}{Hạn chế của Automated Optimization Prompt}
  \begin{itemize}\tightlist
    \item Tốn nhiều lượt chạy (Rollout) mới tìm được prompt tốt.\\
          $\rightarrow$ Chi phí về Rollout Budget
    \item Prompt bị phình to (Các pp sẽ overfit bằng cách gom nhiều ví dụ few-shot vào
          để tăng điểm)\\
          $\rightarrow$ Chi phí về Inference Token Budget.
  \end{itemize}
  \vspace{0.4em}
  Phương pháp GEPA sẽ giải quyết 2 vấn đề trên.
\end{frame}

% ------------------------------------------------------------------ 14
\begin{frame}{1.3\quad Hạn chế của từng phương pháp}{Hạn chế của PEFT}
  \begin{itemize}\tightlist
    \item Dù PEFT đã cắt giảm mạnh chi phí VRAM so với Full Fine-Tuning và giữ được chi
          phí suy luận thấp nhưng:
      \begin{itemize}\tightlist
        \item Chi phí thích nghi lặp lại theo từng tác vụ: Bản chất của PEFT truyền
              thống vẫn là huấn luyện chuyên biệt cho từng bài toán.
        \item Zero-shot Adaption: Khác với Prompting có thể chuyển sang tác vụ mới ngay
              lập tức chỉ bằng vài dòng mô tả, PEFT không thể phản ứng tức thời khi chưa
              có dữ liệu và thời gian train.
      \end{itemize}
  \end{itemize}
\end{frame}

% ------------------------------------------------------------------ 15
\begin{frame}{1.3\quad Hạn chế của từng phương pháp}
  \justifying
  Làm sao để giữ được ưu thế cập nhật trọng số và tốc độ suy luận của LoRA, nhưng xóa bỏ
  hoàn toàn chi phí dữ liệu và thời gian huấn luyện lại mỗi khi gặp tác vụ mới?

  \vspace{0.6em}
  \begin{beamercolorbox}[wd=\dimexpr\linewidth-1em\relax,sep=0.6em]{block title alerted}
    $\rightarrow$ Text-to-LoRA giải quyết nút thắt này bằng kiến trúc Hypernetwork, cho
    phép sinh trực tiếp ma trận trọng số LoRA ngay từ văn bản mô tả tác vụ chỉ trong
    1 lượt chạy xuôi.
  \end{beamercolorbox}
\end{frame}
